\documentclass[11pt]{article}
\usepackage[margin=1in]{geometry}
\usepackage{amsmath,amssymb,amsthm,mathtools}
\usepackage[T1]{fontenc}
\usepackage[utf8]{inputenc}
\usepackage{enumitem}
\usepackage{booktabs,longtable,array}
\usepackage{hyperref}
\usepackage{xcolor}
\usepackage{microtype}
\hypersetup{colorlinks=true, linkcolor=blue!50!black, citecolor=blue!50!black, urlcolor=blue!50!black}

\newtheorem{theorem}{Theorem}[section]
\newtheorem{proposition}[theorem]{Proposition}
\newtheorem{lemma}[theorem]{Lemma}
\newtheorem{corollary}[theorem]{Corollary}
\newtheorem{claim}[theorem]{Claim}
\theoremstyle{definition}
\newtheorem{definition}[theorem]{Definition}
\newtheorem{assumption}[theorem]{Assumption}
\theoremstyle{remark}
\newtheorem{remark}[theorem]{Remark}

\newcommand{\KL}{\mathrm{KL}}
\newcommand{\Ent}{\mathrm{Ent}}
\newcommand{\Prob}{\mathsf{P}}
\newcommand{\E}{\mathbb{E}}
\newcommand{\R}{\mathbb{R}}
\newcommand{\N}{\mathbb{N}}
\newcommand{\Pcal}{\mathcal{P}}
\newcommand{\Mcal}{\mathcal{M}}
\newcommand{\Qcal}{\mathcal{Q}}
\newcommand{\Zcal}{\mathcal{Z}}
\newcommand{\Fcal}{\mathcal{F}}
\newcommand{\Gcal}{\mathcal{G}}
\newcommand{\W}{\mathcal{W}}
\newcommand{\one}{\mathbf{1}}
\newcommand{\dd}{\,d}
\newcommand{\supp}{\operatorname{supp}}
\newcommand{\argmin}{\operatorname*{arg\,min}}
\newcommand{\argmax}{\operatorname*{arg\,max}}
\newcommand{\esssup}{\operatorname*{ess\,sup}}
\newcommand{\ri}{\operatorname{ri}}
\newcommand{\dom}{\operatorname{dom}}
\newcommand{\Law}{\operatorname{Law}}
\newcommand{\Var}{\operatorname{Var}}
\newcommand{\Cov}{\operatorname{Cov}}
\newcommand{\diag}{\operatorname{diag}}
\newcommand{\Id}{\mathrm{Id}}
\newcommand{\Lip}{\operatorname{Lip}}
\newcommand{\osc}{\operatorname{osc}}
\newcommand{\diam}{\operatorname{diam}}
\newcommand{\proj}{\operatorname{proj}}
\newcommand{\e}{\mathrm{e}}

\newenvironment{argument}{\paragraph{Argument.}\begin{enumerate}[leftmargin=2em]}{\end{enumerate}}
\newenvironment{proofoutline}{\paragraph{Proof architecture.}\begin{enumerate}[leftmargin=2em]}{\end{enumerate}}


% Source-faithfulness apparatus used by the paper reconstructions.
\newcommand{\sourceanchor}[1]{\par\smallskip\noindent\textbf{Source anchor.} #1\par\smallskip}
\newcommand{\sourcecomment}[1]{\begin{remark}#1\end{remark}}
\newenvironment{sourceguide}
  {\begin{description}[style=nextline,leftmargin=3.2cm,labelwidth=3cm]}
  {\end{description}}

% Backward-compatible environments used by some of the older reconstructions.
\newenvironment{distilled}[1][]{\begin{theorem}[#1]}{\end{theorem}}
\newenvironment{distillation}{\begin{enumerate}[leftmargin=2em]}{\end{enumerate}}

\newenvironment{commentarynotes}{\begin{remark}\begin{enumerate}[leftmargin=2em]}{\end{enumerate}\end{remark}}

\title{Diffusion Schr\"odinger Bridge Matching}
\author{}
\date{Shi, De Bortoli, Campbell, and Doucet (2023)}
\begin{document}
\maketitle

\section*{Source map and scope}
\begin{sourceguide}
\item[Primary source] Yuyang Shi, Valentin De Bortoli, Andrew Campbell, and Arnaud Doucet, ``Diffusion Schr\"odinger Bridge Matching,'' NeurIPS 2023, arXiv:2303.16852 \cite{ShiDeBortoliCampbellDoucet2023}.
\item[Coverage] Sections 2--4 and Appendix C: bridge matching, Markovian and reciprocal projections, Propositions 2, 4, 5, 7, 9, and 10, Lemma 6, Theorem 8, and the population regression realization of DSBM.
\item[Source order] Equations (1)--(14), Definitions 1 and 3, Propositions 2--10, and Appendix C.1--C.9.
\item[Supplemental references] The argument explicitly relies on Doob transforms, entropic Girsanov, Schr\"odinger-bridge structure, and weak lower semicontinuity; these sources are identified where they enter rather than being absorbed into the paper's proof.  The principal dependencies are \cite{PalmowskiRolski2002,Leonard2012Girsanov,Leonard2014Survey,LeonardRoellyZambrini2014,vanErvenHarremoes2014}.
\end{sourceguide}

This reconstruction follows the source's proof route.  The phrase ``under mild assumptions'' in the main text is replaced by the assumptions actually listed in Appendix C whenever the source gives them.  Steps that the source delegates to cited theorems remain delegated but are stated precisely enough to show the dependency.

\section{Path measures, bridge mixtures, and the Schr\"odinger problem}
Let
\[
 \mathcal C=C([0,T],\R^d),
 \qquad
 \Pcal(\mathcal C)=\{\text{probability laws on }\mathcal C\}.
\]
The reference law \(Q\) is the Markov diffusion
\begin{equation}
 dX_t=f_t(X_t)dt+\sigma_t dB_t,
 \qquad X_0\sim Q_0,
 \tag{1}\label{eq:DSBM-Q}
\end{equation}
with scalar \(\sigma_t>0\).  Write \(Q_{s,t}\) for a two-time marginal, \(Q_{t\mid s}\) for a conditional transition law, and \(Q^{0,T}(\cdot\mid x_0,x_T)\) for the pinned bridge.

For an endpoint coupling \(\Pi_{0,T}\), the bridge mixture
\[
 \Pi=\Pi_{0,T}Q^{0,T}
\]
means
\[
 \Pi(A)=\int Q^{0,T}(A\mid x_0,x_T)\,\Pi_{0,T}(dx_0,dx_T).
\]
Under the Doob $h$-transform assumptions used by the paper \cite{RogersWilliams2000,PalmowskiRolski2002}, the bridge from \(x_0\) to \(x_T\) solves
\begin{equation}
 dX_t^{0,T}
 =\left[f_t(X_t^{0,T})+\sigma_t^2\nabla\log q_{T\mid t}(x_T\mid X_t^{0,T})\right]dt
 +\sigma_t dB_t.
 \tag{2}\label{eq:DSBM-bridge}
\end{equation}

The Schr\"odinger bridge with endpoint laws \(\pi_0,\pi_T\) is
\begin{equation}
 P_{\mathrm{SB}}
 =\argmin\{\KL(P\|Q):P_0=\pi_0,\ P_T=\pi_T\}.
 \tag{6}\label{eq:DSBM-SB}
\end{equation}
By the entropy chain rule with respect to \((X_0,X_T)\), its conditional bridges equal those of \(Q\), so
\[
 P_{\mathrm{SB}}=P^{\mathrm{SB}}_{0,T}Q^{0,T}
\]
and the endpoint coupling solves the static Schr\"odinger problem.

\section{Bridge matching before iteration}
If \(\Pi=\Pi_{0,T}Q^{0,T}\), its bridge drift in \eqref{eq:DSBM-bridge} depends on \(X_T\), hence the mixture is generally not Markov.  Bridge Matching replaces it by a Markov diffusion with the same one-time marginals.  The candidate correction drift is
\begin{equation}
 v_t^*(x)
 =\sigma_t^2\E_\Pi\!
 \left[\nabla\log q_{T\mid t}(X_T\mid X_t)\mid X_t=x\right].
 \tag{3-target}\label{eq:DSBM-vstar}
\end{equation}
At population level it is the minimizer of
\begin{equation}
 \int_0^T\E_\Pi\left[
 \frac{\|\sigma_t^2\nabla\log q_{T\mid t}(X_T\mid X_t)-v_t(X_t)\|^2}{\sigma_t^2}
 \right]dt.
 \tag{3}\label{eq:DSBM-regression}
\end{equation}
The Markovian projection result below explains both the conditional expectation and preservation of all one-time marginals.

\section{The reciprocal class and the two projections}
Let \(\Mcal\) be the class of path laws associated with SDEs
\[
 dX_t=[f_t(X_t)+v_t(X_t)]dt+\sigma_t dB_t
\]
for locally Lipschitz \(v\), in the regularity class needed for the entropy form of Girsanov's theorem.

\begin{definition}[Definition 3]
The reciprocal class of \(Q\) is
\[
 \mathcal R(Q)=\{\Pi_{0,T}Q^{0,T}:\Pi_{0,T}\in\Pcal(\R^d\times\R^d)\}.
\]
Its elements have exactly the same conditional bridges as \(Q\).
\end{definition}

The two maps alternated by IMF are
\[
 \proj_{\Mcal}(\Pi)=\text{the Markovian projection},
 \qquad
 \proj_{\mathcal R(Q)}(P)=P_{0,T}Q^{0,T}.
\]
They optimize opposite orientations of relative entropy.

\section{Appendix assumptions for the Markovian projection}
\sourceanchor{Appendix C.1, assumptions A1--A3 and Lemma 11.}

The paper's main text says ``under mild assumptions.''  Appendix C uses the following concrete package.

\begin{assumption}[A1: well-posedness and growth]
The functions \(f\), \(\sigma\), and
\[
 (t,x)\longmapsto
 \E_{\Pi_{T\mid t}}[\nabla\log q_{T\mid t}(X_T\mid X_t)\mid X_t=x]
\]
are locally Lipschitz.  There are \(C>0\) and continuous \(\psi\ge0\) such that
\[
 \|f_t(x)\|\le C(1+\|x\|),
 \qquad
 C^{-1}\le\sigma_t\le C,
\]
and
\[
 \left\|\E[\nabla\log q_{T\mid t}(X_T\mid X_t)\mid X_t=x]\right\|
 \le C\psi(t)(1+\|x\|).
\]
\end{assumption}

\begin{assumption}[A2: a good Doob function]
For every \(x_0\), \(\Pi_{T\mid0}(\cdot\mid x_0)\ll Q_{T\mid0}(\cdot\mid x_0)\), with bounded density
\[
 \phi_{T\mid0}(x_T\mid x_0)
 =\frac{d\Pi_{T\mid0}}{dQ_{T\mid0}}(x_T\mid x_0).
\]
Define
\begin{equation}
 \phi_{t\mid0}(x_t\mid x_0)
 =\int\phi_{T\mid0}(x_T\mid x_0)Q_{T\mid t}(dx_T\mid x_t).
 \tag{32}\label{eq:DSBM-phi}
\end{equation}
For every \(x_0\), both \(1/\phi_{t\mid0}\) and \(A\phi_{t\mid0}\) are bounded, where
\begin{equation}
 Ah=\partial_th+\langle f_t,\nabla h\rangle+\frac{\sigma_t^2}{2}\Delta h.
 \tag{31}\label{eq:DSBM-generator}
\end{equation}
\end{assumption}

Under A2, the cited Doob-transform results give
\begin{equation}
 \phi_{\cdot\mid0}\in C^{1,2},
 \qquad
 A\phi_{\cdot\mid0}=0.
 \tag{Lemma 11}\label{eq:DSBM-backward-kolm}
\end{equation}

\begin{assumption}[A3: gradient growth]
For every \(x_0\), there is \(C_{x_0}\) such that
\[
 \|\nabla\log\phi_{t\mid0}(x_t\mid x_0)\|
 \le C_{x_0}(1+\|x_0\|+\|x_t\|).
\]
\end{assumption}

These assumptions are sufficient for the Doob transform, Girsanov entropy identity, and uniqueness of the Fokker--Planck equation used in Proposition 2.  The source notes that they are not intended to be minimal.

\section{Proposition 2: reverse-KL Markovian projection}
\sourceanchor{Section 3.1, Proposition 2; Appendix C.1, equations (31)--(35).}

\begin{proposition}[Proposition 2]
Assume A1--A3 and let \(\Pi=\Pi_{0,T}Q^{0,T}\).  Define \(M^*=\proj_{\Mcal}(\Pi)\) by
\begin{equation}
 dX_t^*=[f_t(X_t^*)+v_t^*(X_t^*)]dt+\sigma_t dB_t,
 \tag{Definition 1}\label{eq:DSBM-Mstar}
\end{equation}
with \(v^*\) from \eqref{eq:DSBM-vstar}.  Then
\[
 M^*=\argmin_{M\in\Mcal}\KL(\Pi\|M),
\]
all one-time marginals are preserved,
\[
 M_t^*=\Pi_t\quad(0\le t\le T),
\]
and
\begin{align}
 \KL(\Pi\|M^*)
 &=\frac12\int_0^T\E_{\Pi_{0,t}}
 \left[
 \frac{\|U_t-\E[U_t\mid X_t]\|^2}{\sigma_t^2}
 \right]dt,
 \tag{P2}\label{eq:DSBM-P2-value}\\
 U_t
 &=\sigma_t^2\E_\Pi[\nabla\log q_{T\mid t}(X_T\mid X_t)\mid X_0,X_t].
\end{align}
\end{proposition}

\subsection{Step 1: identify the semimartingale drift of the bridge mixture}
Conditionally on \(X_0=x_0\), the path law is the Doob transform of \(Q_{\mid0}\) by \(\phi_{T\mid0}\).  Equations \eqref{eq:DSBM-generator} and \eqref{eq:DSBM-backward-kolm} give the transformed generator
\[
 \frac{A(h\phi)-hA\phi}{\phi}
 =Ah+\sigma_t^2\langle\nabla h,\nabla\log\phi\rangle.
\]
Hence, conditionally on \(X_0\), \(\Pi\) has drift
\begin{equation}
 f_t(X_t)+\sigma_t^2\nabla\log\phi_{t\mid0}(X_t\mid X_0).
 \tag{33}\label{eq:DSBM-conditional-drift}
\end{equation}

Bayes' identity gives
\[
 q_{t\mid0}(x_t\mid x_0)\phi_{t\mid0}(x_t\mid x_0)
 =\pi_{t\mid0}(x_t\mid x_0).
\]
Using
\[
 q_{T\mid0}(x_T\mid x_0)q_{t\mid0,T}(x_t\mid x_0,x_T)
 =q_{t\mid0}(x_t\mid x_0)q_{T\mid t}(x_T\mid x_t)
\]
in \eqref{eq:DSBM-phi}, differentiate with respect to \(x_t\) to obtain
\begin{equation}
 \nabla\log\phi_{t\mid0}(x_t\mid x_0)
 =\E_\Pi[\nabla\log q_{T\mid t}(X_T\mid X_t)
          \mid X_0=x_0,X_t=x_t].
 \tag{34--35}\label{eq:DSBM-phi-score}
\end{equation}
Thus \(U_t\) is exactly the nonreference part of the semimartingale drift of \(\Pi\).

\subsection{Step 2: minimize the Girsanov energy}
For a Markov candidate
\[
 dX_t=[f_t(X_t)+v_t(X_t)]dt+\sigma_t dB_t,
\]
the entropic Girsanov theorem \cite{Leonard2012Girsanov} gives
\begin{equation}
 \KL(\Pi\|M)
 =\frac12\int_0^T\E_\Pi
 \left[\frac{\|U_t-v_t(X_t)\|^2}{\sigma_t^2}\right]dt.
 \tag{G}\label{eq:DSBM-girsanov-objective}
\end{equation}
For each \(t\), conditional expectation is the orthogonal projection from \(L^2(\Pi)\) onto the subspace of \(X_t\)-measurable vectors.  Therefore
\[
 v_t^*(X_t)=\E[U_t\mid X_t]
\]
uniquely minimizes the integrand, and the residual energy is \eqref{eq:DSBM-P2-value}.

\subsection{Step 3: preservation of marginals}
The Fokker--Planck weak equation for the possibly non-Markov law \(\Pi\) contains the conditional current drift
\[
 \E[f_t(X_t)+U_t\mid X_t]
 =f_t(X_t)+v_t^*(X_t).
\]
This is precisely the drift of \(M^*\).  Since the two laws have the same initial marginal and solve the same Fokker--Planck equation, uniqueness of probability solutions of the Fokker--Planck equation \cite{BogachevKrasovitskiiShaposhnikov2021} under A1 and A3 gives \(M_t^*=\Pi_t\) for every \(t\).

\section{Proposition 4: forward-KL reciprocal projection}
\sourceanchor{Section 3.1, Proposition 4; Appendix C.2.}

\begin{proposition}[Proposition 4]
For any path law \(P\),
\[
 \proj_{\mathcal R(Q)}(P)=P_{0,T}Q^{0,T}
 =\argmin_{\Pi\in\mathcal R(Q)}\KL(P\|\Pi).
\]
\end{proposition}

Disintegrate both laws with respect to the endpoints.  The chain rule gives
\begin{equation}
 \KL(P\|\Pi)
 =\KL(P_{0,T}\|\Pi_{0,T})
 +\int\KL(P^{0,T}_{x_0,x_T}\|Q^{0,T}_{x_0,x_T})P_{0,T}(dx_0,dx_T).
 \tag{C.2}\label{eq:DSBM-chain-endpoint}
\end{equation}
The second term is independent of \(\Pi_{0,T}\).  The first is minimized, with value zero, by \(\Pi_{0,T}=P_{0,T}\).

\section{Proposition 5: the Markov--reciprocal intersection}
\sourceanchor{Section 3.1, Proposition 5; Appendix C.3.}

The source invokes L\'eonard's existence and uniqueness theorem \cite{Leonard2014Survey,LeonardRoellyZambrini2014} rather than reproving the Schr\"odinger system.  The assumptions quoted in Appendix C.3 are:
\begin{enumerate}[leftmargin=2em]
\item \(Q\) is Markov with equal endpoint marginals \(Q_0=Q_T=\bar Q\);
\item the endpoint density of \(Q\) has a lower bound \(dQ_{0,T}/d(\bar Q\otimes\bar Q)\ge e^{-A(x_0)-A(x_T)}\);
\item a complementary exponential integrability condition involving a measurable \(B\ge0\);
\item an irreducibility condition at some intermediate time \(t_0\);
\item finite endpoint entropies and integrability of \(A+B\) under \(\pi_0,\pi_T\).
\end{enumerate}
Under these conditions there is a unique Schr\"odinger bridge.  Moreover:

\begin{proposition}[Proposition 5]
If \(P\in\Mcal\cap\mathcal R(Q)\), \(P_0=\pi_0\), \(P_T=\pi_T\), and \(\KL(P\|Q)<\infty\), then
\[
 P=P_{\mathrm{SB}}.
\]
\end{proposition}

The reason is the source theorem's characterization of the unique minimizer as the unique finite-entropy law that is both Markov and in the reciprocal class with the required endpoints.  This proposition is the fixed-point identification used by IMF.

\section{Lemma 6: the two KL Pythagorean identities}
\sourceanchor{Section 3.2, Lemma 6; Appendix C.4, equations (36)--(37).}

\begin{lemma}[Lemma 6]
For \(M\in\Mcal\) and \(\Pi\in\mathcal R(Q)\), whenever the displayed entropies are finite,
\begin{align}
 \KL(\Pi\|M)
 &=\KL(\Pi\|\proj_{\Mcal}\Pi)
  +\KL(\proj_{\Mcal}\Pi\|M),
 \tag{36}\label{eq:DSBM-Pyth-M}\\
 \KL(M\|\Pi)
 &=\KL(M\|\proj_{\mathcal R(Q)}M)
  +\KL(\proj_{\mathcal R(Q)}M\|\Pi).
 \tag{37}\label{eq:DSBM-Pyth-R}
\end{align}
\end{lemma}

\subsection{Markovian identity}
Let \(U_t\) be the conditional bridge-mixture drift from Proposition 2, and let
\[
 \bar U_t=\E[U_t\mid X_t].
\]
The three Girsanov formulas are
\begin{align*}
 2\KL(\Pi\|M)
 &=\int\E\|U_t-v_t(X_t)\|^2\sigma_t^{-2}dt,\\
 2\KL(\Pi\|\proj_{\Mcal}\Pi)
 &=\int\E\|U_t-\bar U_t\|^2\sigma_t^{-2}dt,\\
 2\KL(\proj_{\Mcal}\Pi\|M)
 &=\int\E\|\bar U_t-v_t(X_t)\|^2\sigma_t^{-2}dt.
\end{align*}
The cross term vanishes because
\[
 \E\langle U_t-\bar U_t,\bar U_t-v_t(X_t)\rangle=0.
\]
This is the conditional-expectation Pythagorean theorem and proves \eqref{eq:DSBM-Pyth-M}.  Appendix C expands the squares instead of naming the Hilbert-space projection.

\subsection{Reciprocal identity}
Set \(\Pi^*=M_{0,T}Q^{0,T}\).  The Radon--Nikodym ratio between \(\Pi^*\) and \(\Pi\) depends only on \((X_0,X_T)\):
\[
 \log\frac{d\Pi^*}{d\Pi}
 =\log\frac{dM_{0,T}}{d\Pi_{0,T}}(X_0,X_T).
\]
Since \(M\) and \(\Pi^*\) share endpoint law \(M_{0,T}\), its expectation is the same under both.  Expanding
\[
 \log\frac{dM}{d\Pi}
 =\log\frac{dM}{d\Pi^*}+\log\frac{d\Pi^*}{d\Pi}
\]
and integrating gives \eqref{eq:DSBM-Pyth-R}.

\begin{remark}[Why convex projection theory is not being used]
Appendix C.9 gives an explicit three-time binary counterexample showing that the set of Markov laws is not convex.  If \(p_1,p_2\) are Markov chains with parameters \(\alpha_i,\beta_i\), their half-mixture fails the Markov identity by
\[
 \Delta=\alpha_1\beta_1(\beta_0-\alpha_0)(\beta_2-\alpha_2).
\]
Thus the first Pythagorean identity is a special quadratic-drift identity, not an application of Csisz\'ar's convex-set projection theorem \cite{Csiszar1975}.
\end{remark}

\section{Iterative Markovian Fitting}
\sourceanchor{Section 3.2, equation (8), Proposition 7, Theorem 8; Appendix C.5--C.6.}

Initialize \(P_0\in\mathcal R(Q)\) with \((P_0)_0=\pi_0\), \((P_0)_T=\pi_T\), and alternate
\begin{equation}
 P_{2n+1}=\proj_{\Mcal}(P_{2n}),
 \qquad
 P_{2n+2}=\proj_{\mathcal R(Q)}(P_{2n+1}).
 \tag{8}\label{eq:DSBM-IMF}
\end{equation}
The Markovian projection preserves every one-time marginal, while the reciprocal projection preserves the endpoint joint law.  Hence every iterate has endpoint marginals \(\pi_0,\pi_T\).

\begin{proposition}[Proposition 7]
Assume Proposition 5 and Lemma 6 apply to every iterate and to \(P_{\mathrm{SB}}\).  Then
\[
 \KL(P_{n+1}\|P_{\mathrm{SB}})
 \le\KL(P_n\|P_{\mathrm{SB}})<\infty,
\]
and
\[
 \KL(P_n\|P_{n+1})\longrightarrow0.
\]
\end{proposition}

For either parity, the bridge belongs to the opposite projection class, so the corresponding identity in Lemma 6 gives
\begin{equation}
 \KL(P_n\|P_{\mathrm{SB}})
 =\KL(P_n\|P_{n+1})
  +\KL(P_{n+1}\|P_{\mathrm{SB}}).
 \tag{tel}\label{eq:DSBM-telescope-one}
\end{equation}
The telescoping device follows the reverse-KL analogue of R\"uschendorf's IPF argument \cite{Ruschendorf1995}.

Iterating,
\begin{equation}
 \KL(P_0\|P_{\mathrm{SB}})
 =\sum_{i=0}^N\KL(P_i\|P_{i+1})
  +\KL(P_{N+1}\|P_{\mathrm{SB}}).
 \tag{C.5}\label{eq:DSBM-telescope}
\end{equation}
The series of nonnegative increments is finite, proving the result.

\begin{theorem}[Theorem 8]
Assume the conditions above, coercivity of \(P\mapsto\KL(P\|P_{\mathrm{SB}})\), weak closedness of the Markov and reciprocal classes along the iterates, and weak lower semicontinuity of relative entropy \cite{vanErvenHarremoes2014}.  Then the IMF sequence has the unique fixed point
\[
 P^*=P_{\mathrm{SB}},
\]
and
\[
 \KL(P_n\|P^*)\to0.
\]
\end{theorem}

\begin{proof}
All iterates lie in the sublevel set
\[
 \mathcal K=\{P:\KL(P\|P_{\mathrm{SB}})
 \le\KL(P_0\|P_{\mathrm{SB}})\}.
\]
Coercivity makes \(\mathcal K\) relatively compact.  Choose a weakly convergent subsequence of the Markov iterates,
\[
 M_{n_j}=P_{2n_j+1}\Rightarrow M^*,
\]
and a further subsequence of the adjacent reciprocal iterates,
\[
 \Pi_{n_{j_k}}=P_{2n_{j_k}+2}\Rightarrow\Pi^*.
\]
Closedness gives \(M^*\in\Mcal\) and \(\Pi^*\in\mathcal R(Q)\).  By lower semicontinuity and Proposition 7,
\[
 0\le\KL(M^*\|\Pi^*)
 \le\liminf_k\KL(M_{n_{j_k}}\|\Pi_{n_{j_k}})=0.
\]
Thus \(M^*=\Pi^*=:P^*\).  Endpoint preservation passes to the limit, so Proposition 5 identifies \(P^*=P_{\mathrm{SB}}\).

Along this subsequence, lower semicontinuity plus equality of the limit is used by the source to conclude that the entropy to \(P^*\) reaches zero; monotonicity from Proposition 7 then forces the full sequence to zero.
\end{proof}

\begin{remark}[Assumptions hidden by ``mild assumptions'']
The source does not prove in general that its SDE-defined Markov class is weakly closed, nor that the reciprocal class is weakly closed for every reference bridge kernel.  It invokes these properties in Appendix C.6.  Coercivity of \(\KL(\cdot\|P_{\mathrm{SB}})\) is standard on a Polish path space because entropy sublevels are tight relative to a fixed probability law, but the passage from weak convergence to convergence of the entropy value requires more than lower semicontinuity alone.  The proof as printed should therefore be read with the stated compactness, closedness, and entropy-convergence conditions included among the theorem's assumptions.
\end{remark}

\section{Proposition 9: forward and backward representations}
\sourceanchor{Section 4, Proposition 9; Appendix C.7, equation (38).}

Time-reversing the pinned bridge \eqref{eq:DSBM-bridge} gives
\begin{equation}
 dY_t^{0,T}
 =\left[-f_{T-t}(Y_t^{0,T})
 +\sigma_{T-t}^2\nabla\log q_{T-t\mid0}(Y_t^{0,T}\mid x_0)\right]dt
 +\sigma_{T-t}dB_t,
 \quad Y_0^{0,T}=x_T.
 \tag{38}\label{eq:DSBM-reverse-bridge}
\end{equation}
Applying the same Markovian-projection argument in forward and reversed time gives:

\begin{proposition}[Proposition 9]
For \(\Pi=\Pi_{0,T}Q^{0,T}\), the same Markovian projection \(M^*\) has the forward representation
\begin{equation}
 dX_t=
 \left[f_t(X_t)+\sigma_t^2
 \E_\Pi[\nabla\log q_{T\mid t}(X_T\mid X_t)\mid X_t]\right]dt
 +\sigma_t dB_t,
 \quad X_0\sim\Pi_0,
 \tag{11}\label{eq:DSBM-forward}
\end{equation}
and the reversed representation
\begin{equation}
 dY_t=
 \left[-f_{T-t}(Y_t)+\sigma_{T-t}^2
 \E_\Pi[\nabla\log q_{T-t\mid0}(Y_t\mid X_0)\mid Y_t]\right]dt
 +\sigma_{T-t}dB_t,
 \quad Y_0\sim\Pi_T.
 \tag{12}\label{eq:DSBM-backward}
\end{equation}
\end{proposition}

The source proof is one paragraph: apply Proposition 2 to the reversed bridge formula \eqref{eq:DSBM-reverse-bridge}.  The two conditional expectations are the forward and backward current drifts of the same Markov law.

\section{Population regressions and Algorithm 1}
The conditional expectations in Proposition 9 are learned by squared-error regression.  For a family \(v_\theta\), the forward objective is
\begin{equation}
 \theta^*=\argmin_\theta
 \int_0^T\E_{\Pi_{t,T}}
 \left[
 \frac{\|\sigma_t^2\nabla\log q_{T\mid t}(X_T\mid X_t)-v_\theta(t,X_t)\|^2}{\sigma_t^2}
 \right]dt.
 \tag{10}\label{eq:DSBM-forward-loss}
\end{equation}
Its unrestricted population minimizer is \(v_t^*\) from \eqref{eq:DSBM-vstar}.  The backward objective is
\begin{equation}
 \varphi^*=\argmin_\varphi
 \int_0^T\E_{\Pi_{0,t}}
 \left[
 \frac{\|\sigma_t^2\nabla\log q_{t\mid0}(X_t\mid X_0)-v_\varphi(t,X_t)\|^2}{\sigma_t^2}
 \right]dt.
 \tag{14}\label{eq:DSBM-backward-loss}
\end{equation}
Its population minimizer is the backward conditional expectation in \eqref{eq:DSBM-backward}.

A reciprocal projection is sampled by first drawing \((X_0,X_T)\) from the current Markov law and then drawing a path from \(Q^{0,T}(\cdot\mid X_0,X_T)\).  A Markovian projection is approximated by fitting one of the regressions and simulating the resulting SDE.  DSBM alternates backward and forward representations so that each fitted SDE is initialized from an exact endpoint law: backward from \(\pi_T\), then forward from \(\pi_0\).  This is the paper's device for preventing endpoint bias from accumulating when the regressions are not exact.

\section{Proposition 10: exact DSBM recovers the IPF sequence}
\sourceanchor{Section 4, Proposition 10; Appendix C.8.}

Let \((\widetilde P_n)\) be the exact path-space IPF sequence initialized at \(Q\).  The cited DSB analysis \cite{DeBortoliThorntonHengDoucet2021} implies each \(\widetilde P_n\) is both Markov and reciprocal with respect to \(Q\):
\[
 \widetilde P_n\in\Mcal\cap\mathcal R(Q).
\]
Assume the regression families are rich enough to realize the exact conditional expectations, and initialize DSBM with \(\Pi^0_{0,T}=Q_{0,T}\).

\begin{proposition}[Proposition 10]
Under exact population regression, the Markov laws generated by DSBM coincide with the IPF iterates:
\[
 M_n=\widetilde P_n
 \qquad(n\ge0).
\]
\end{proposition}

\begin{proof}
The claim holds at initialization.  Assume \(M_{2n+1}=\widetilde P_{2n+1}\).  Since the IPF iterate already belongs to \(\mathcal R(Q)\),
\[
 \Pi_{2n+1}
 =\proj_{\mathcal R(Q)}(M_{2n+1})
 =M_{2n+1}.
\]
Since it is also Markov,
\[
 \proj_{\Mcal}(\Pi_{2n+1})=M_{2n+1}.
\]
The next forward DSBM law is initialized at \(\pi_0\) and uses the conditional transition law of this exact projection.  The IPF update is the same operation, written as time reversal followed by reinitialization:
\[
 \widetilde P_{2n+2}=\pi_0\widetilde P_{2n+1\mid0}.
\]
Hence \(M_{2n+2}=\widetilde P_{2n+2}\).  The backward step is identical with the roles of \(0\) and \(T\) exchanged.  Induction proves the result.
\end{proof}

\section{Brownian bridges and the deterministic limit}
For \(f=0\) and constant \(\sigma\),
\[
 X_t^{0,T}=\left(1-\frac tT\right)x_0+\frac tT x_T
 +\sigma\left(B_t-\frac tTB_T\right),
\]
and
\[
 dX_t^{0,T}=\frac{x_T-X_t^{0,T}}{T-t}dt+\sigma dB_t.
\]
The forward target becomes
\[
 \frac{X_T-X_t}{T-t},
\]
so the Bridge Matching loss is
\[
 \E\left\|\frac{X_T-X_t}{T-t}-v_\theta(t,X_t)\right\|^2.
\]
As \(\sigma\downarrow0\), the bridge collapses to linear interpolation and the regression becomes the Flow Matching velocity regression.  This limiting observation motivates the algorithm, but the projection proofs above assume \(\sigma_t>0\).

\section{Proof dependencies and source-level cautions}
\begin{enumerate}[leftmargin=2em]
\item Proposition 2 is not proved from elementary SDE facts alone.  It depends on the Doob \(h\)-transform results cited to Palmowski--Rolski, the entropic Girsanov identity cited to L\'eonard, and uniqueness of the Fokker--Planck equation.
\item Proposition 5 is a direct use of L\'eonard's Schr\"odinger-bridge existence and Markov--reciprocal characterization theorem.
\item Lemma 6 is exact once the entropy identities are available; its Markovian half is an \(L^2\) orthogonal-decomposition calculation and its reciprocal half is the endpoint entropy chain rule.
\item Theorem 8 uses compactness of entropy sublevels, weak closedness of both alternating classes, and a convergence step for the entropy value.  These are compressed into ``mild assumptions'' in the main text and are not established in full generality by the paper.
\item Proposition 9 is stated with ``under mild conditions'' and its appendix proof says only that it is analogous to Proposition 2; the reverse-time Doob/Girsanov conditions must therefore be supplied as well.
\item Proposition 10 is an exact-population statement.  The practical algorithm has approximation, time-discretization, and finite-sample errors, so equality with IPF does not survive automatically outside that idealization.
\end{enumerate}

\begin{thebibliography}{99}
\bibitem{ChenGeorgiouPavon2021} Yongxin Chen, Tryphon T. Georgiou, and Michele Pavon. \emph{Stochastic Control Liaisons: Richard Sinkhorn Meets Gaspard Monge on a Schr\"odinger Bridge}. SIAM Review, 63(2):249--313, 2021. doi:10.1137/20M1339982. arXiv:2005.10963.
\bibitem{Leonard2014Survey} Christian L\'eonard. \emph{A Survey of the Schr\"odinger Problem and Some of Its Connections with Optimal Transport}. Discrete and Continuous Dynamical Systems--A, 34(4):1533--1574, 2014. doi:10.3934/dcds.2014.34.1533. arXiv:1308.0215.
\bibitem{Fortet1940} Robert Fortet. \emph{Resolution d'un systeme d'equations de M. Schrodinger}. Journal de mathematiques pures et appliquees, 9e serie, 19(1--4):83--105, 1940.
\bibitem{SinkhornKnopp1967} Richard Sinkhorn and Paul Knopp. \emph{Concerning Nonnegative Matrices and Doubly Stochastic Matrices}. Pacific Journal of Mathematics, 21(2):343--348, 1967.
\bibitem{FranklinLorenz1989} Joel Franklin and Jens Lorenz. \emph{On the Scaling of Multidimensional Matrices}. Linear Algebra and its Applications, 114/115:717--735, 1989. doi:10.1016/0024-3795(89)90490-4.
\bibitem{Carroll2004} Joseph E. Carroll. \emph{Birkhoff's Contraction Coefficient}. Linear Algebra and its Applications, 389:227--234, 2004. doi:10.1016/j.laa.2004.02.039.
\bibitem{EvesonNussbaum1995} Simon P. Eveson and Roger D. Nussbaum. \emph{An Elementary Proof of the Birkhoff--Hopf Theorem}. Mathematical Proceedings of the Cambridge Philosophical Society, 117(1):31--55, 1995.
\bibitem{BlanchetMurthy2019} Jose Blanchet and Karthyek R. A. Murthy. \emph{Quantifying Distributional Model Risk via Optimal Transport}. Mathematics of Operations Research, 44(2):565--600, 2019. doi:10.1287/moor.2018.0936. arXiv:1604.01446.
\bibitem{GaoKleywegt2023} Rui Gao and Anton J. Kleywegt. \emph{Distributionally Robust Stochastic Optimization with Wasserstein Distance}. Mathematics of Operations Research, 48(2):603--655, 2023. doi:10.1287/moor.2022.1275. arXiv:1604.02199.
\bibitem{MohajerinEsfahaniKuhn2018} Peyman Mohajerin Esfahani and Daniel Kuhn. \emph{Data-Driven Distributionally Robust Optimization Using the Wasserstein Metric: Performance Guarantees and Tractable Reformulations}. Mathematical Programming, 171(1--2):115--166, 2018. doi:10.1007/s10107-017-1172-1. arXiv:1505.05116.
\bibitem{WangGaoXie2025} Jie Wang, Rui Gao, and Yao Xie. \emph{Sinkhorn Distributionally Robust Optimization}. Operations Research, 2025. doi:10.1287/opre.2023.0294. arXiv:2109.11926.
\bibitem{Girsanov1960} I. V. Girsanov. \emph{On Transforming a Certain Class of Stochastic Processes by Absolutely Continuous Substitution of Measures}. Theory of Probability and Its Applications, 5(3):285--301, 1960. doi:10.1137/1105027.
\bibitem{CameronMartin1945} R. H. Cameron and W. T. Martin. \emph{Transformations of Wiener Integrals under a General Class of Linear Transformations}. Transactions of the American Mathematical Society, 58(2):184--219, 1945.
\bibitem{Yeh1978} J. Yeh. \emph{Transformation of Conditional Wiener Integrals under Translation and the Cameron--Martin Translation Theorem}. Tohoku Mathematical Journal, 30(4):505--515, 1978.
\bibitem{Kakutani1948} Shizuo Kakutani. \emph{On Equivalence of Infinite Product Measures}. Annals of Mathematics, 49(1):214--224, 1948.
\bibitem{ShiDeBortoliCampbellDoucet2023} Yuyang Shi, Valentin De Bortoli, Andrew Campbell, and Arnaud Doucet. \emph{Diffusion Schr\"odinger Bridge Matching}. Advances in Neural Information Processing Systems, 2023. arXiv:2303.16852.
\bibitem{LiuLipmanNickelKarrerTheodorouChen2024} Guan-Horng Liu, Yaron Lipman, Maximilian Nickel, Brian Karrer, Evangelos A. Theodorou, and Ricky T. Q. Chen. \emph{Generalized Schr\"odinger Bridge Matching}. International Conference on Learning Representations, 2024. arXiv:2310.02233.
\bibitem{DomingoEnrichDrozdzalKarrerChen2025} Carles Domingo-Enrich, Michal Drozdzal, Brian Karrer, and Ricky T. Q. Chen. \emph{Adjoint Matching: Fine-Tuning Flow and Diffusion Generative Models with Memoryless Stochastic Optimal Control}. arXiv preprint arXiv:2409.08861, 2025.

\bibitem{Birkhoff1957} Garrett Birkhoff. \emph{Extensions of Jentzsch's Theorem}. Transactions of the American Mathematical Society, 85(1):219--227, 1957.
\bibitem{Sinkhorn1964} Richard Sinkhorn. \emph{A Relationship Between Arbitrary Positive Matrices and Doubly Stochastic Matrices}. Annals of Mathematical Statistics, 35(2):876--879, 1964.
\bibitem{Sinkhorn1967Prescribed} Richard Sinkhorn. \emph{Diagonal Equivalence to Matrices with Prescribed Row and Column Sums}. American Mathematical Monthly, 74(4):402--405, 1967.
\bibitem{Bauer1965} Friedrich L. Bauer. \emph{An Elementary Proof of the Hopf Inequality for Positive Operators}. Numerische Mathematik, 7:331--337, 1965.
\bibitem{Bushell1973} P. J. Bushell. \emph{Hilbert's Metric and Positive Contraction Mappings in a Banach Space}. Archive for Rational Mechanics and Analysis, 52:330--338, 1973.

\end{thebibliography}

\end{document}
